%% ============================================================
%%  CopilotScope Quality Measurement Framework
%%  A multi-signal, telemetry-driven model for assessing
%%  AI coding-assistant session quality.
%%
%%  Compile: pdflatex quality_measurement_framework.tex  (twice)
%%  Requires: amsmath, amssymb, booktabs, hyperref, xcolor,
%%            geometry, graphicx, algorithm2e (or algorithmicx)
%% ============================================================
\documentclass[11pt,a4paper]{article}

\usepackage[T1]{fontenc}
\usepackage[utf8]{inputenc}
\usepackage{lmodern}
\usepackage[margin=2.5cm]{geometry}
\usepackage{amsmath,amssymb,amsthm}
\usepackage{booktabs}
\usepackage{array}
\usepackage{xcolor}
\usepackage{graphicx}
\usepackage{hyperref}
\usepackage{microtype}
\usepackage{parskip}
\usepackage[ruled,vlined,linesnumbered]{algorithm2e}

\hypersetup{
  colorlinks=true,
  linkcolor=blue!70!black,
  citecolor=green!50!black,
  urlcolor=blue!60!black,
  pdftitle={CopilotScope Quality Measurement Framework},
  pdfauthor={CopilotScope Project},
}

\newtheorem{definition}{Definition}
\newtheorem{proposition}{Proposition}
\newtheorem{remark}{Remark}

\newcommand{\R}{\mathbb{R}}
\newcommand{\N}{\mathbb{N}}
\newcommand{\calI}{\mathcal{I}}
\newcommand{\calW}{\mathcal{W}}
\newcommand{\clamp}[3]{\operatorname{clamp}\!\left(#1,\,#2,\,#3\right)}

% ---- title block ----
\title{%
  \textbf{CopilotScope Quality Measurement Framework}\\[0.4em]
  \large A Multi-Signal, Telemetry-Driven Model for Assessing\\
  AI Coding-Assistant Session Quality
}
\author{CopilotScope Project}
\date{July 2026 \quad (draft)}

% ============================================================
\begin{document}
\maketitle
\tableofcontents
\newpage

% ------------------------------------------------------------
\section{Introduction}
\label{sec:intro}
% ------------------------------------------------------------

AI coding assistants (GitHub Copilot, Claude Code, Cursor) generate significant
telemetry during development sessions: token counts, acceptance/rejection events,
time-to-first-token (TTFT) measurements, error traces, and user messages.
This telemetry is rich but heterogeneous---different signals arrive at different
times, from different editor surfaces, and at different granularities.

\textbf{CopilotScope} is an OpenTelemetry-based collector and dashboard that
ingests this telemetry in real time and derives actionable quality scores.
This paper formally defines the eight algorithms that constitute the quality
measurement pipeline, examines their mathematical properties, and provides
numerical examples for each.

The pipeline design follows three principles:
\begin{enumerate}
  \item \textbf{Telemetry-only}---no LLM judge, no prompt inspection, no
        network round-trip at scoring time.
  \item \textbf{Explainability by construction}---every score decomposes into
        named, human-readable sub-scores with reasons.
  \item \textbf{Confidence-aware}---a score without data does not silently
        collapse to a false neutral; it carries a confidence of~0.
\end{enumerate}

% ------------------------------------------------------------
\section{System Architecture Overview}
\label{sec:architecture}
% ------------------------------------------------------------

Telemetry flows from the editor (VS~Code, JetBrains, CLI) via OTLP/HTTP or
OTLP/proto to the \texttt{CopilotScope.Collector} service, which:

\begin{enumerate}
  \item Decodes payloads using a zero-dependency custom protobuf parser
        (\S\ref{sec:proto}).
  \item Routes spans and metrics to the appropriate \texttt{CopilotSession}
        aggregate.
  \item Runs the \texttt{InsightPipeline} on every ingest, producing a
        \texttt{QualityReport} and up to six \texttt{InsightReport}s.
\end{enumerate}

The quality pipeline consists of the algorithms defined in
Sections~\ref{sec:quality-engine}--\ref{sec:percentile}.

% ------------------------------------------------------------
\section{Composite Session Quality Engine}
\label{sec:quality-engine}
% ------------------------------------------------------------

\subsection{Component Definitions}

The engine aggregates up to six independent components.  Let
$\calI \subseteq \{1,\dots,6\}$ denote the set of \emph{informative}
components---those with at least one observed data point
($n_i > 0$).

\begin{table}[h!]
\centering
\caption{Quality components, base weights, and sub-formulas.}
\label{tab:components}
\begin{tabular}{@{}llll@{}}
\toprule
Component & Symbol & Base weight $w_i$ & Value $v_i$ \\
\midrule
Reliability & $R$    & 0.25 & $(1 - \hat{e})^2$ \\
Acceptance  & $A$    & 0.20 & $0.6\alpha + 0.4\bar{s}$ \\
Friction    & $\Phi$ & 0.20 & $\bar{\tau}$ (TFRA mean turn score) \\
Latency     & $\Lambda$ & 0.15 & $L(p_{50})$ (log-linear curve) \\
Feedback    & $F$    & 0.10 & $u^+ / (u^+ + u^-)$ \\
Efficiency  & $E$    & 0.10 & $\overline{\mathrm{eff}}$ \\
\bottomrule
\end{tabular}
\end{table}

\paragraph{Reliability.}
LLM (chat) errors are weighted twice as heavily as tool errors,
reflecting their greater impact on session flow.
\[
  \hat{e} = \frac{2\,e_{\text{LLM}} + e_{\text{tool}}}
                 {\max(1,\;2\,c_{\text{LLM}} + c_{\text{tool}})},
  \qquad
  v_R = \clamp{(1 - \hat{e})^2}{0}{1}.
\]
The quadratic penalty means that a~10\% error rate reduces reliability by
nearly 19\% relative to zero errors, while a~50\% error rate yields
$v_R = 0.25$---ensuring discrimination at high error rates.

\paragraph{Acceptance.}
Let $a$ = edits accepted, $j$ = edits rejected,
$\bar{s}$ = mean survival score (see \S\ref{sec:edit-survival}).
\[
  \alpha = \frac{a}{a+j}, \qquad
  v_A = \clamp{0.6\alpha + 0.4\bar{s}}{0}{1}.
\]

\paragraph{Friction.}
$\bar{\tau} = \frac{1}{N}\sum_{i=1}^N \tau_i$, where $\tau_i$ is the
per-turn TFRA score (see \S\ref{sec:tfra}).

\paragraph{Latency.}
$v_\Lambda = L(p_{50})$ where $L$ is defined in \S\ref{sec:latency-utility}.

\paragraph{Feedback.}
$v_F = u^+ / (u^+ + u^-)$, the thumbs-up ratio from explicit user votes.

\paragraph{Efficiency.}
The efficiency score is the mean of the available sub-signals:
\[
  v_E = \frac{1}{|\mathcal{P}|}\sum_{p \in \mathcal{P}} e_p,
\]
where $\mathcal{P}$ includes (i)~the cache-hit ratio
$r_c = C / (I + C)$ and (ii)~the turns-per-invocation score
\[
  e_{\mathrm{tpi}} = \begin{cases} 1 & \text{if } \bar{t} \leq 8 \\
    \clamp{1 - (\bar{t}-8)/17}{0}{1} & \text{otherwise,} \end{cases}
\]
with $\bar{t}$ = mean turns per agent invocation.  The turn-per-invocation
penalty saturates at 25 turns ($e_{\mathrm{tpi}} = 0$).

\subsection{Composite Score with Weight Renormalisation}
\label{sec:composite}

The key design decision of v2 is to \emph{renormalise} over informative
components only, eliminating the v1 bug where absent components contributed
a neutral prior at full weight, compressing the score toward~79--80 for
any ordinary session.

\begin{definition}[Coverage]
  $C = \sum_{i \in \calI} w_i \in (0, 1]$.
\end{definition}

\begin{definition}[Composite quality score]
\[
  \boxed{Q = \frac{100}{C} \sum_{i \in \calI} w_i\, v_i}
\]
When $\calI = \varnothing$ the system returns the neutral prior $Q = 70$
at confidence~0.
\end{definition}

\begin{proposition}
$Q \in [0, 100]$ for all $v_i \in [0,1]$.
\end{proposition}
\begin{proof}
  $Q = (100/C)\sum_{i \in \calI} w_i v_i \leq (100/C)\sum_{i \in \calI} w_i = 100$.
  The lower bound follows analogously.
\end{proof}

\subsection{Confidence}

The \emph{confidence} signal communicates how much data backs the score.
Each component contributes a \emph{sample ramp}
$\rho_i = \min(1,\, n_i/5)$ that reaches 1.0 after five samples.

\begin{definition}[Confidence]
\[
  \rho = \frac{\sum_{i \in \calI} w_i\,\rho_i}{C},
  \qquad
  \boxed{\text{Confidence} = C \cdot \rho \in [0,1].}
\]
\end{definition}

A confidence of 1.0 implies both full weight coverage ($C=1$) and at
least five samples per component.  The interpretation ``score 91 at
confidence 0.2'' reads as ``early but promising'', not ``certain''.

\subsection{Relative Score}

Given a history window of $H$ recent session scores
$\{Q_1,\dots,Q_H\}$ with mean $\mu$ and standard deviation $\sigma$:
\[
  z = \frac{Q - \mu}{\sigma}, \qquad
  \text{PercentileRank} = \frac{|\{j : Q_j \leq Q\}|}{H}.
\]

\begin{table}[h!]
\centering
\caption{Grade thresholds and relative labels.}
\begin{tabular}{@{}lll@{}}
\toprule
Score $Q$ & Grade & Relative label (by percentile) \\
\midrule
$\geq 85$ & excellent & above baseline ($\geq 0.75$) \\
$\geq 70$ & good & at baseline ($[0.35, 0.75)$) \\
$\geq 55$ & fair & below baseline ($< 0.35$) \\
$\geq 40$ & poor & \\
$< 40$    & critical & \\
\bottomrule
\end{tabular}
\end{table}

\subsection{Worked Example}
\label{sec:quality-engine-example}

We score one realistic session end-to-end.  Five of the six components are
informative; the user cast no explicit thumbs votes, so Feedback has
$n_F = 0$ and drops out of $\calI$.  The measured inputs are:
$e_{\text{LLM}} = 1$ error over $c_{\text{LLM}} = 12$ chat calls and
$e_{\text{tool}} = 2$ errors over $c_{\text{tool}} = 30$ tool calls;
$a = 14$ accepted and $j = 6$ rejected edits with mean survival
$\bar{s} = 0.6936$; mean TFRA turn score $\bar{\tau} = 0.87$ over
$N = 10$ turns; median TTFT $p_{50} = 1{,}950\,\text{ms}$ over 6 latency
samples; and efficiency sub-signals $r_c = 0.7714$ with
$\bar{t} = 5$ turns per invocation over 2 agent invocations.

\paragraph{Step 1 --- Reliability $v_R$.}
We form the weighted error rate, with chat errors counted twice, then
apply the quadratic penalty.
\[
  \hat{e} = \frac{2 \times 1 + 2}{\max(1,\; 2 \times 12 + 30)}
          = \frac{4}{54}
          = 0.0741,
\]
\[
  v_R = \clamp{(1 - 0.0741)^2}{0}{1}
      = 0.9259^2
      = 0.8573.
\]

\paragraph{Step 2 --- Acceptance $v_A$.}
The acceptance ratio is blended with the survival mean, 60/40.
\[
  \alpha = \frac{14}{14+6} = 0.70,
  \qquad
  v_A = \clamp{0.6 \times 0.70 + 0.4 \times 0.6936}{0}{1}
      = 0.4200 + 0.2774
      = 0.6974.
\]

\paragraph{Step 3 --- Friction $v_\Phi$ and Latency $v_\Lambda$.}
Friction passes the TFRA mean through unchanged; latency pushes the median
TTFT through the log-linear curve of \S\ref{sec:latency-utility}.
\[
  v_\Phi = \bar{\tau} = 0.8700,
\]
\[
  v_\Lambda = L(1950)
            = 1 - \frac{\ln(1950/300)}{\ln(10000/300)}
            = 1 - \frac{1.8718}{3.5066}
            = 1 - 0.5338
            = 0.4662.
\]

\paragraph{Step 4 --- Efficiency $v_E$.}
Both sub-signals are available: the cache ratio, and the
turns-per-invocation score, which is 1 because $\bar{t} = 5 \leq 8$.
\[
  v_E = \frac{r_c + e_{\mathrm{tpi}}}{2}
      = \frac{0.7714 + 1}{2}
      = 0.8857.
\]

\paragraph{Step 5 --- Coverage $C$.}
Feedback is uninformative, so only the five observed base weights sum
into the coverage.
\[
  C = 0.25 + 0.20 + 0.20 + 0.15 + 0.10 = 0.90.
\]

\paragraph{Step 6 --- Composite score $Q$.}
The five component values are weighted, summed, and renormalised by the
coverage so the missing Feedback weight does not drag the score down.
\[
  \sum_{i \in \calI} w_i v_i
  = 0.25 \times 0.8573
  + 0.20 \times 0.6974
  + 0.20 \times 0.8700
  + 0.15 \times 0.4662
  + 0.10 \times 0.8857
\]
\[
  \phantom{\sum_{i \in \calI} w_i v_i}
  = 0.2143 + 0.1395 + 0.1740 + 0.0699 + 0.0886
  = 0.6863,
\]
\[
  Q = \frac{100}{C} \sum_{i \in \calI} w_i v_i
    = \frac{100}{0.90} \times 0.6863
    = 76.26 \quad (\text{grade: good}).
\]

\paragraph{Step 7 --- Confidence.}
Each component's sample ramp saturates at five samples; only Efficiency,
with 2 invocations, is still ramping.
\[
  \rho_R = \rho_A = \rho_\Phi = \rho_\Lambda = 1,
  \qquad
  \rho_E = \min(1,\; 2/5) = 0.40,
\]
\[
  \rho = \frac{0.25 + 0.20 + 0.20 + 0.15 + 0.10 \times 0.40}{0.90}
       = \frac{0.8400}{0.90}
       = 0.9333,
\]
\[
  \text{Confidence} = C \cdot \rho = 0.90 \times 0.9333 = 0.84.
\]

\paragraph{Step 8 --- Relative score.}
Against a history window of $H = 20$ recent sessions with mean
$\mu = 72.4$ and standard deviation $\sigma = 6.8$, of which 14 scored at
or below $Q$:
\[
  z = \frac{76.26 - 72.4}{6.8} = \frac{3.86}{6.8} = 0.57,
  \qquad
  \text{PercentileRank} = \frac{14}{20} = 0.70
  \quad (\text{at baseline}).
\]

The session reads as ``good (76.26) at confidence 0.84, at baseline for
this user'': a well-evidenced, slightly above-average session whose main
drag is latency ($v_\Lambda = 0.4662$).

\begin{table}[h!]
\centering
\caption{Summary of the worked composite-score example.}
\label{tab:quality-engine-worked}
\begin{tabular}{@{}llll@{}}
\toprule
Symbol & Meaning & Value (with unit) & What happens at this step \\
\midrule
$\hat{e}$ & weighted error rate & 0.0741 (7.4\,\%) & $4/54$; chat errors count double \\
$v_R$ & reliability value & 0.8573 & quadratic penalty on $1-\hat{e}$ \\
$v_A$ & acceptance value & 0.6974 & $0.6\alpha + 0.4\bar{s}$ \\
$v_\Phi$ & friction value & 0.8700 & TFRA mean passed through \\
$v_\Lambda$ & latency value & 0.4662 & $L(p_{50})$ at 1{,}950 ms \\
$v_E$ & efficiency value & 0.8857 & mean of $r_c$ and $e_{\mathrm{tpi}}$ \\
$C$ & coverage & 0.90 & Feedback ($w=0.10$) uninformative \\
$Q$ & composite score & 76.26 points & weighted sum renormalised by $C$ \\
Confidence & data backing & 0.84 & Efficiency ramp still at 0.40 \\
$z$, PercentileRank & relative position & 0.57, 0.70 & at baseline in $H=20$ window \\
\bottomrule
\end{tabular}
\end{table}

\subsection{Numerical Example}

\begin{table}[h!]
\centering
\caption{Example composite scores for representative sessions.}
\label{tab:quality-examples}
\begin{tabular}{@{}lrrl@{}}
\toprule
Scenario & $Q$ & Confidence & Grade \\
\midrule
Perfect (all components, no errors) & 94.2 & 1.00 & excellent \\
Two LLM errors, reliability only    & 64.0 & 0.25 & good \\
No data at all                      & 70.0 & 0.00 & good (prior) \\
High rejection rate, good latency   & 61.5 & 0.52 & fair \\
\bottomrule
\end{tabular}
\end{table}

% ------------------------------------------------------------
\section{Turn-level Friction \& Repair Analysis (TFRA)}
\label{sec:tfra}
% ------------------------------------------------------------

Each agent-invocation trace is a \emph{turn}.  The TFRA model scores every
turn on $[0,1]$ starting at 1.0 and applying penalty deductions for
observable friction signals.

\subsection{Per-turn Score}

\begin{definition}[TFRA turn score]
\[
  \tau_i = \max\!\left(0,\;
    1
    - 0.35 \cdot \min(e^{\text{LLM}}_i, 2)
    - 0.15 \cdot \min(e^{\text{tool}}_i, 3)
    - \Delta^{\text{lat}}_i
    - \Delta^{\text{rep}}_i
  \right)
\]
\end{definition}

\paragraph{Latency penalty $\Delta^{\text{lat}}_i$ (self-referential).}
Let $m$ = session median TTFT.  A turn is penalised only when its
latency is anomalous \emph{relative to its own session}---so a
slow model is not penalised for being slow, only for being inconsistent.
\[
  \Delta^{\text{lat}}_i =
  \begin{cases}
    0.30 & \text{if } \text{TTFT}_i \geq 3m \\
    0.15 & \text{if } 1.5m \leq \text{TTFT}_i < 3m \\
    0    & \text{otherwise.}
  \end{cases}
\]

\paragraph{Repair-loop penalty $\Delta^{\text{rep}}_i$.}
A repair loop is the trace signature of an agent retrying out of a failure:
unusually many tool calls per chat call \emph{and} tool errors present.
Let $m_{\text{tr}} = \text{Median}\!\left(c^{\text{tool}}_j / \max(1, c^{\text{LLM}}_j)\right)$.
\[
  \Delta^{\text{rep}}_i = 0.10 \cdot \mathbf{1}\!\left[
    e^{\text{tool}}_i > 0
    \;\wedge\;
    \frac{c^{\text{tool}}_i}{\max(1, c^{\text{LLM}}_i)} \geq 2\, m_{\text{tr}}
    \;\wedge\;
    c^{\text{tool}}_i \geq 3
  \right].
\]

\subsection{Session Median}

\begin{definition}[Median]
For a sorted list $v_{(1)} \leq \cdots \leq v_{(n)}$:
\[
  \text{Median}(V) =
  \begin{cases}
    v_{(\lfloor n/2 \rfloor + 1)} & n \text{ odd} \\
    \dfrac{v_{(n/2)} + v_{(n/2+1)}}{2} & n \text{ even.}
  \end{cases}
\]
\end{definition}

\subsection{Best/Worst Turn Selection}

\[
  i^* = \arg\max_i \tau_i \quad (\text{ties: } \arg\min_i \text{TTFT}_i),
  \qquad
  i_- = \arg\min_i \tau_i \quad (\text{ties: } \arg\max_i \text{TTFT}_i).
\]

\subsection{Penalty Budget Summary}

\begin{table}[h!]
\centering
\caption{Maximum penalties per friction source.}
\begin{tabular}{@{}lll@{}}
\toprule
Source & Max penalty & Capping behaviour \\
\midrule
LLM errors & $0.70$ & capped at 2 errors \\
Tool errors & $0.45$ & capped at 3 errors \\
Latency outlier & $0.30$ & one penalty only (worst tier wins) \\
Repair loop & $0.10$ & binary \\
\midrule
Total possible & $1.55$ & but $\tau_i = \max(0, \cdot)$ floors at 0 \\
\bottomrule
\end{tabular}
\end{table}

\subsection{Worked Example}
\label{sec:tfra-example}

We score one realistic 5-turn session whose per-turn trace signals are
given below.

\begin{table}[h!]
\centering
\caption{Per-turn trace signals for the worked TFRA session.}
\label{tab:tfra-inputs}
\begin{tabular}{@{}crrrrrr@{}}
\toprule
Turn $i$ & $\text{TTFT}_i$ & $e^{\text{LLM}}_i$ & $e^{\text{tool}}_i$
  & $c^{\text{LLM}}_i$ & $c^{\text{tool}}_i$
  & $c^{\text{tool}}_i / \max(1, c^{\text{LLM}}_i)$ \\
\midrule
1 & $900\,\text{ms}$   & 0 & 0 & 1 & 2 & 2 \\
2 & $1{,}000\,\text{ms}$ & 1 & 0 & 1 & 1 & 1 \\
3 & $1{,}200\,\text{ms}$ & 0 & 1 & 2 & 4 & 2 \\
4 & $1{,}500\,\text{ms}$ & 0 & 0 & 1 & 2 & 2 \\
5 & $3{,}800\,\text{ms}$ & 0 & 2 & 2 & 8 & 4 \\
\bottomrule
\end{tabular}
\end{table}

\paragraph{Step 1 --- Session medians.}
Both penalty rules are self-referential, so we first fix the session's
median TTFT and median tool ratio ($n = 5$, odd, so the middle value).
\[
  m = \text{Median}(900, 1000, 1200, 1500, 3800) = 1{,}200\,\text{ms},
\]
\[
  m_{\text{tr}} = \text{Median}(2, 1, 2, 2, 4) = 2.
\]

\paragraph{Step 2 --- Penalty thresholds.}
The medians fix the concrete cutoffs every turn is judged against.
\[
  1.5\,m = 1{,}800\,\text{ms}, \qquad
  3\,m = 3{,}600\,\text{ms}, \qquad
  2\,m_{\text{tr}} = 4.
\]

\paragraph{Step 3 --- Worst turn written in full ($\tau_5$).}
Turn 5 trips all three non-LLM penalties: two tool errors cost
$0.15 \times \min(2,3) = 0.30$; its TTFT of $3{,}800\,\text{ms} \geq
3{,}600\,\text{ms}$ costs $\Delta^{\text{lat}}_5 = 0.30$; and the repair
test fires because $e^{\text{tool}}_5 = 2 > 0$, the tool ratio
$8/2 = 4 \geq 2\,m_{\text{tr}} = 4$, and $c^{\text{tool}}_5 = 8 \geq 3$,
costing $\Delta^{\text{rep}}_5 = 0.10$.
\[
  \tau_5 = \max\!\left(0,\;
    1 - 0.35 \times \min(0, 2) - 0.15 \times \min(2, 3) - 0.30 - 0.10
  \right)
\]
\[
  \phantom{\tau_5}
  = \max(0,\; 1 - 0 - 0.30 - 0.30 - 0.10)
  = \max(0,\; 0.30)
  = 0.30.
\]

\paragraph{Step 4 --- Remaining turns.}
The other four turns each trip at most one penalty: all TTFTs are below
$1{,}800\,\text{ms}$, and turn~3 fails the repair test since its tool
ratio $2 < 4$.
\[
  \tau_1 = 1.00, \qquad
  \tau_2 = 1 - 0.35 \times \min(1,2) = 0.65, \qquad
  \tau_3 = 1 - 0.15 \times \min(1,3) = 0.85, \qquad
  \tau_4 = 1.00.
\]

\paragraph{Step 5 --- Session mean $\bar{\tau}$.}
The five turn scores are averaged; this mean is what feeds the Friction
component of \S\ref{sec:quality-engine}.
\[
  \bar{\tau} = \frac{1.00 + 0.65 + 0.85 + 1.00 + 0.30}{5}
             = \frac{3.80}{5}
             = 0.76.
\]

\paragraph{Step 6 --- Best/worst turn selection.}
Turns 1 and 4 tie at $\tau = 1.00$, so the tie-break picks the smaller
TTFT; the worst turn is unambiguous.
\[
  i^* = 1 \quad (900\,\text{ms} < 1{,}500\,\text{ms}),
  \qquad
  i_- = 5 \quad (\tau_5 = 0.30).
\]

\begin{table}[h!]
\centering
\caption{Summary of the worked TFRA example.}
\label{tab:tfra-worked}
\begin{tabular}{@{}llll@{}}
\toprule
Symbol & Meaning & Value (with unit) & What happens at this step \\
\midrule
$m$ & session median TTFT & 1{,}200 ms & middle of 5 sorted TTFTs \\
$m_{\text{tr}}$ & median tool ratio & 2 & middle of 5 tool ratios \\
$\Delta^{\text{lat}}_5$ & latency penalty, turn 5 & 0.30 & $3{,}800 \geq 3m = 3{,}600$ ms \\
$\Delta^{\text{rep}}_5$ & repair penalty, turn 5 & 0.10 & errors + ratio $4 \geq 4$ + calls $\geq 3$ \\
$\tau_5$ & worst turn score & 0.30 & three penalties stack \\
$\tau_1, \ldots, \tau_4$ & other turn scores & 1.00, 0.65, 0.85, 1.00 & at most one penalty each \\
$\bar{\tau}$ & session mean & 0.76 & $3.80 / 5$; feeds Friction \\
$i^*$, $i_-$ & best/worst turns & 1, 5 & tie at 1.00 broken by TTFT \\
\bottomrule
\end{tabular}
\end{table}

% ------------------------------------------------------------
\section{Edit Survival Analysis}
\label{sec:edit-survival}
% ------------------------------------------------------------

\subsection{Signals}

Two editor-side signals measure how much of generated code survives after
the observation window closes:

\begin{itemize}
  \item \textbf{Four-gram overlap} $F \in [0,1]$: fraction of 4-gram character
        shingles from generated code still present verbatim.
  \item \textbf{No-revert} $R \in \{0,1\}$: whether the edit was never undone,
        averaged over edits to yield $\bar{R} \in [0,1]$.
\end{itemize}

\subsection{Combined Survival Score}

\begin{definition}[Edit survival]
When both signals are present:
\[
  \boxed{S = 0.4\,\bar{F} + 0.6\,\bar{R}.}
\]
When only one signal is available, $S = \bar{F}$ or $S = \bar{R}$ respectively.
\end{definition}

\begin{remark}
The no-revert weight (0.6) exceeds the four-gram weight (0.4) because code
can be legitimately \emph{refactored}---low four-gram overlap but no revert---while
the semantic intent of the AI suggestion survives.  No-revert is the stronger
durability signal.
\end{remark}

\begin{table}[h!]
\centering
\caption{Interpretation thresholds.}
\begin{tabular}{@{}ll@{}}
\toprule
$S$ & Interpretation \\
\midrule
$\geq 0.8$ & Durable --- vast majority survives untouched \\
$[0.5, 0.8)$ & Partial rework --- suggestions land but get meaningfully edited \\
$< 0.5$ & Heavy rewrite --- most generated code is replaced or reverted \\
\bottomrule
\end{tabular}
\end{table}

% ---- 4.3 Four-gram shingling mechanics ----
\subsection{Four-gram Shingling --- Exact Mechanics}
\label{sec:edit-survival-4gram}

Given a generated code string $g$ and its surviving form $s$, the
\emph{4-gram set} is defined as:
\[
  \mathcal{G}_4(t) = \bigl\{\, t[i : i+4] \;\bigm|\; 0 \leq i \leq |t| - 4 \,\bigr\}.
\]
Using sets (not multisets) prevents double-counting repeated patterns
(e.g.\ indentation runs \verb|"    "|) and keeps $F \in [0,1]$ regardless of
code length.

\begin{definition}[Four-gram overlap for a single edit]
\[
  F = \frac{|\mathcal{G}_4(g) \cap \mathcal{G}_4(s)|}{|\mathcal{G}_4(g)|}.
\]
The per-session mean $\bar{F} = \tfrac{1}{n}\sum_{k=1}^{n} F_k$ averages over
all $n$ edits in the session.
\end{definition}

\paragraph{Mini-example: single-character change.}
Consider a generated line \texttt{"\ \ \ \ return x + 1"} (16 characters)
that the user changes to \texttt{"\ \ \ \ return x + 2"}.

\begin{itemize}
  \item $|\mathcal{G}_4(g)| = 13$ (4-gram positions $0\ldots 12$).
  \item The only differing 4-gram is at position 12:
        generated \texttt{" + 1"} $\neq$ surviving \texttt{" + 2"}.
  \item $|\mathcal{G}_4(g) \cap \mathcal{G}_4(s)| = 12$.
  \item $F = 12/13 \approx 0.923$\,---high survival despite a user edit.
\end{itemize}

\paragraph{Mini-example: semantic rewrite.}
Generated: \texttt{"if x > 0:\textbackslash n\ \ \ \ return True"} (22 chars, 19 grams).
Surviving: \texttt{"if x >= 0:\textbackslash n\ \ \ \ return x != 0"} (25 chars, 22 grams).

The shared 4-grams are structural tokens such as \texttt{"if x"}, \texttt{"\ \ \ \ "} and
\texttt{"retu"}\,---\,11 of 19 generated grams survive (see Table~\ref{tab:4gram-rewrite}).

\begin{table}[h!]
\centering
\caption{4-gram overlap for the semantic-rewrite mini-example.}
\label{tab:4gram-rewrite}
\begin{tabular}{@{}lrl@{}}
\toprule
Metric & Value & Explanation \\
\midrule
$|\mathcal{G}_4(g)|$ & 19 & 4-grams in generated code \\
$|\mathcal{G}_4(s)|$ & 22 & 4-grams in surviving code \\
$|\mathcal{G}_4(g) \cap \mathcal{G}_4(s)|$ & 11 & shared grams (structure/keywords) \\
$|\mathcal{G}_4(g) \setminus \mathcal{G}_4(s)|$ & 8 & lost grams (operator, value tokens) \\
$F$ & $11/19 \approx 0.579$ & substantial rewrite detected \\
\bottomrule
\end{tabular}
\end{table}

% ---- 4.4 Worked numerical example ----
\subsection{Worked Example --- Full Session (4 Edits)}
\label{sec:edit-survival-example}

We trace a realistic 4-edit Python session end-to-end.  For each edit,
$F_k$ is computed via four-gram overlap and $R_k \in \{0,1\}$ records
whether the edit was ever undone.

\begin{table}[h!]
\centering
\caption{Per-edit signal values for the worked session.}
\label{tab:edit-per-edit}
\begin{tabular}{@{}clllrrl@{}}
\toprule
$k$ & Generated (excerpt) & Surviving (excerpt) & Event & $F_k$ & $R_k$ & Note \\
\midrule
1 & \texttt{def add(a,b): return a+b} & identical & kept verbatim
    & $1.000$ & $1$ & all grams match \\
2 & \texttt{return x * y} & \texttt{return x * y * 2}
    & append & $0.857$ & $1$ & last gram changed \\
3 & \texttt{if x > 0: return True} & \texttt{if x >= 0: return x != 0}
    & rewrite & $0.579$ & $1$ & logic replaced \\
4 & \texttt{\# TODO: implement later} & (empty) & reverted
    & $0.000$ & $0$ & undo triggered \\
\bottomrule
\end{tabular}
\end{table}

\paragraph{Step 1 — Session means.}
We average the four per-edit similarity scores and, separately, the four
revert indicators, giving the session's two aggregate survival signals.
\[
  \bar{F} = \frac{F_1 + F_2 + F_3 + F_4}{4}
           = \frac{1.000 + 0.857 + 0.579 + 0.000}{4}
           = \frac{2.436}{4} = 0.6090.
\]
\[
  \bar{R} = \frac{R_1 + R_2 + R_3 + R_4}{4}
           = \frac{1 + 1 + 1 + 0}{4}
           = \frac{3}{4} = 0.7500.
\]

\paragraph{Step 2 --- Combined survival score.}
The two means are blended with the fixed weights: 0.4 on textual
survival and 0.6 on the sturdier revert signal.
\[
  S = 0.4\,\bar{F} + 0.6\,\bar{R}
    = 0.4 \times 0.6090 + 0.6 \times 0.7500
    = 0.2436 + 0.4500
    = 0.6936.
\]

\paragraph{Step 3 --- Interpretation.}
$S = 0.6936 \in [0.5, 0.8)$ $\Rightarrow$ \textbf{partial rework}.
Despite one full revert and one substantial rewrite, two clean edits and
three no-revert outcomes keep the session above the heavy-rewrite threshold.

\begin{remark}
The four-gram signal is sensitive to small changes: edit~2 loses only one
4-gram out of seven ($F_2 = 6/7 \approx 0.857$), confirming that appending
\texttt{* 2} is a minor extension rather than a replacement.
\end{remark}

\begin{table}[h!]
\centering
\caption{Summary of the worked edit-survival example.}
\label{tab:edit-survival-worked}
\begin{tabular}{@{}llll@{}}
\toprule
Symbol & Meaning & Value (with unit) & What happens at this step \\
\midrule
$F_1, \ldots, F_4$ & per-edit four-gram survival & 1.000, 0.857, 0.579, 0.000 & shingle overlap per edit \\
$R_1, \ldots, R_4$ & per-edit revert indicators & 1, 1, 1, 0 & edit 4 was undone \\
$\bar{F}$ & mean four-gram survival & 0.6090 & $2.436 / 4$ \\
$\bar{R}$ & mean revert survival & 0.7500 & $3 / 4$ \\
$S$ & combined survival score & 0.6936 & $0.4\bar{F} + 0.6\bar{R}$ \\
Band & interpretation & $[0.5, 0.8)$ & partial rework \\
\bottomrule
\end{tabular}
\end{table}

% ---- 4.5 Weight sensitivity ----
\subsection{Weight Sensitivity}
\label{sec:edit-survival-sensitivity}

Let $w_F + w_R = 1$, $w_F, w_R \geq 0$.  Table~\ref{tab:weight-sensitivity}
evaluates five $(w_F, w_R)$ pairs on the four scenarios from
Table~\ref{tab:edit-per-edit} (using the session-level means) plus the
\emph{refactored-but-kept} scenario $(\bar{F}=0.25,\bar{R}=0.95)$.

\begin{table}[h!]
\centering
\caption{Combined survival score $S$ under alternative weights.
         $\star$ = CopilotScope default.  Cells show $S$ (interpretation:
         \textbf{D} = durable, \textbf{P} = partial rework, \textbf{H} = heavy rewrite).}
\label{tab:weight-sensitivity}
\begin{tabular}{@{}lrrrrrrr@{}}
\toprule
Scenario & $\bar{F}$ & $\bar{R}$ &
  $w_F{=}0.2$ & $w_F{=}0.3$ & $w_F{=}0.4^\star$ &
  $w_F{=}0.5$ & $w_F{=}0.6$ \\
\midrule
Durable (verbatim)
  & 0.890 & 0.975
  & 0.958\,D & 0.949\,D & 0.941\,D & 0.933\,D & 0.924\,D \\
Refactored (semantic keep)
  & 0.250 & 0.950
  & 0.810\,D & 0.740\,P & 0.670\,P & 0.600\,P & 0.530\,P \\
Worked session (§\ref{sec:edit-survival-example})
  & 0.609 & 0.750
  & 0.572\,P & 0.608\,P & 0.694\,P & 0.680\,P & 0.665\,P \\
Mixed signals
  & 0.575 & 0.650
  & 0.605\,P & 0.628\,P & 0.620\,P & 0.613\,P & 0.605\,P \\
Mostly reverted
  & 0.275 & 0.150
  & 0.145\,H & 0.187\,H & 0.200\,H & 0.213\,H & 0.225\,H \\
\bottomrule
\end{tabular}
\end{table}

The critical row is \emph{Refactored (semantic keep)}:
with the default weights $(0.4, 0.6)$ the score is 0.670 (partial rework,
correctly capturing that AI intent survived).  A four-gram-dominant
weighting $w_F{=}0.6$ drops it to 0.530---still partial rework but
closer to the heavy-rewrite boundary, understating the durability of a
semantically accepted but syntactically refactored suggestion.

% ---- 4.6 Connection to acceptance component ----
\subsection{Connection to the Acceptance Component}
\label{sec:edit-survival-acceptance}

The per-session survival score $\bar{S}$ feeds directly into the
\textbf{Acceptance} component of the Composite Quality Engine
(Table~\ref{tab:components}):
\[
  v_A = \clamp{0.6\,\alpha + 0.4\,\bar{S}}{0}{1},
\]
where $\alpha = a/(a+j)$ is the edit-acceptance ratio.

This two-level design captures two distinct moments in the editing workflow:

\begin{itemize}
  \item $\alpha$ records the \emph{immediate} accept/reject decision at
        suggestion time.
  \item $\bar{S}$ records the \emph{long-run} durability after the
        observation window closes.
\end{itemize}

Table~\ref{tab:acceptance-examples} illustrates how $v_A$ discriminates
between scenarios that look identical on $\alpha$ alone.

\begin{table}[h!]
\centering
\caption{Acceptance component $v_A$ for representative scenarios.}
\label{tab:acceptance-examples}
\begin{tabular}{@{}lrrrrl@{}}
\toprule
Scenario & $\alpha$ & $\bar{S}$ & $0.6\alpha$ & $0.4\bar{S}$ & $v_A$ \\
\midrule
High accept + high survival
  & 0.90 & 0.85 & 0.540 & 0.340 & 0.880 \\
High accept + heavy rewrite
  & 0.90 & 0.25 & 0.540 & 0.100 & 0.640 \\
Moderate accept + high survival
  & 0.60 & 0.90 & 0.360 & 0.360 & 0.720 \\
Low accept + durable keeps
  & 0.40 & 0.95 & 0.240 & 0.380 & 0.620 \\
Low accept + low survival
  & 0.40 & 0.20 & 0.240 & 0.080 & 0.320 \\
\bottomrule
\end{tabular}
\end{table}

\begin{remark}
``High accept + heavy rewrite'' scores $v_A = 0.640$, \emph{lower} than
``Moderate accept + high survival'' ($v_A = 0.720$).  A user who clicks
\emph{Accept} but immediately rewrites the code is not a high-quality
interaction; the survival term correctly penalises this pattern.
\end{remark}

% ------------------------------------------------------------
\section{Acceptance-Weighted Throughput}
\label{sec:throughput}
% ------------------------------------------------------------

\subsection{Component Metrics}

Let $a$ = edits accepted, $j$ = edits rejected, $L$ = lines added,
$T$ = number of turns ($T \geq 1$), $O$ = output tokens.

\[
  \alpha = \frac{a}{a+j}, \quad
  \ell_{\text{adj}} = \frac{L \cdot \alpha}{T}, \quad
  \ell_{\text{tok}} = \frac{L}{O/1000} \quad (\text{optional}).
\]

\subsection{Throughput Score}

\begin{definition}[Throughput]
\[
  \boxed{\text{Score} = \clamp{0.6\alpha + 0.4\min\!\left(1,\;\frac{\ell_{\text{adj}}}{20}\right)}{0}{1}.}
\]
\end{definition}

The LOC component saturates at $\ell_{\text{adj}} = 20$ accepted LOC per turn.
Acceptance ratio dominates (weight 0.6) because high throughput with many
rejections is churn, not productivity.

\begin{table}[h!]
\centering
\caption{Numerical examples for throughput scenarios.}
\begin{tabular}{@{}lrrrrr@{}}
\toprule
Scenario & $\alpha$ & $\ell_{\text{adj}}$ & Score & Finding \\
\midrule
High acceptance + high LOC & 0.89 & 17.8 & 0.89 & high throughput \\
Low acceptance rate & 0.40 & 4.0 & 0.32 & miss the mark \\
Good acceptance, low LOC & 0.83 & 2.5 & 0.55 & moderate \\
Perfect, moderate LOC & 1.00 & 8.0 & 0.76 & high throughput \\
\bottomrule
\end{tabular}
\end{table}

\subsection{Worked Example}
\label{sec:throughput-example}

We trace one realistic refactoring session end-to-end.  The assistant
proposed $a + j = 20$ edits over the session: $a = 14$ edits were accepted
and $j = 6$ edits were rejected.  The accepted edits contributed
$L = 220$ lines added across $T = 10$ turns, and the model emitted
$O = 5{,}500$ output tokens in total.

\paragraph{Step 1 --- Acceptance ratio $\alpha$.}
First we measure what fraction of the assistant's proposals the user
actually kept, which is the quality signal that dominates the score.
\[
  \alpha = \frac{a}{a+j}
         = \frac{14}{14+6}
         = \frac{14}{20}
         = 0.70 \quad (\text{dimensionless, i.e.\ } 70\,\%).
\]

\paragraph{Step 2 --- Acceptance-adjusted LOC per turn $\ell_{\text{adj}}$.}
Next we convert raw volume into \emph{useful} volume: the lines added are
discounted by the acceptance ratio, then normalised per turn so long
sessions are not favoured.
\[
  \ell_{\text{adj}} = \frac{L \cdot \alpha}{T}
                    = \frac{220\,\text{LOC} \times 0.70}{10\,\text{turns}}
                    = \frac{154\,\text{LOC}}{10\,\text{turns}}
                    = 15.4\ \text{accepted LOC/turn}.
\]

\paragraph{Step 3 --- Token-normalised density $\ell_{\text{tok}}$ (optional).}
The optional diagnostic relates lines produced to tokens spent; it does not
enter the score but flags verbose models.
\[
  \ell_{\text{tok}} = \frac{L}{O/1000}
                    = \frac{220\,\text{LOC}}{5{,}500/1000\,\text{kTok}}
                    = \frac{220\,\text{LOC}}{5.5\,\text{kTok}}
                    = 40\ \text{LOC/kTok}.
\]

\paragraph{Step 4 --- LOC saturation term.}
The LOC component is capped so that sheer volume beyond 20 accepted
LOC/turn earns no extra credit; here the session sits just below the cap.
\[
  \min\!\left(1,\;\frac{\ell_{\text{adj}}}{20}\right)
  = \min\!\left(1,\;\frac{15.4}{20}\right)
  = \min(1,\;0.77)
  = 0.77.
\]

\paragraph{Step 5 --- Weighted combination.}
We blend the two components with the fixed weights: 0.6 on acceptance
(quality) and 0.4 on saturated LOC rate (volume).
\[
  0.6\alpha + 0.4 \times 0.77
  = 0.6 \times 0.70 + 0.4 \times 0.77
  = 0.420 + 0.308
  = 0.728.
\]

\paragraph{Step 6 --- Clamp.}
Finally the result is clamped to $[0,1]$; since $0.728$ already lies inside
the interval, the clamp is the identity here.
\[
  \text{Score} = \clamp{0.728}{0}{1} = 0.728.
\]

The session scores $0.728$ --- between the ``Good acceptance, low LOC''
and ``Perfect, moderate LOC'' scenarios above: a solid, near-saturation
volume of accepted code, held back mainly by the 30\,\% rejection rate.

\begin{table}[h!]
\centering
\caption{Summary of the worked throughput example.}
\label{tab:throughput-worked}
\begin{tabular}{@{}llll@{}}
\toprule
Symbol & Meaning & Value (with unit) & What happens at this step \\
\midrule
$a$ & edits accepted & 14 edits & raw input from telemetry \\
$j$ & edits rejected & 6 edits & raw input from telemetry \\
$\alpha$ & acceptance ratio & 0.70 (70\,\%) & $14/20$; quality signal \\
$L$ & lines added & 220 LOC & raw input from telemetry \\
$T$ & number of turns & 10 turns & normalises volume per turn \\
$\ell_{\text{adj}}$ & accepted LOC per turn & 15.4 LOC/turn & $220 \times 0.70 / 10$ \\
$O$ & output tokens & 5{,}500 tokens & raw input from telemetry \\
$\ell_{\text{tok}}$ & LOC per kilotoken & 40 LOC/kTok & $220/5.5$; diagnostic only \\
--- & saturation term & 0.77 & $\min(1, 15.4/20)$; below cap \\
Score & throughput score & 0.728 & $0.6 \times 0.70 + 0.4 \times 0.77$, clamped \\
\bottomrule
\end{tabular}
\end{table}

% ------------------------------------------------------------
\section{Latency-Utility Model}
\label{sec:latency-utility}
% ------------------------------------------------------------

\subsection{Motivation}

A single p50 latency figure does not capture user-perceived impact.  A session
with ten fast turns and two very slow ones has a low p50 but meaningful
abandonment risk.  The latency-utility model applies a curve to \emph{every
sample}, capturing distributional shape.

\subsection{Utility Function}

The curve is log-linear in the interval $[300\,\text{ms}, 10{,}000\,\text{ms}]$,
reflecting the multiplicative perception of latency (Weber--Fechner law applied
to time):

\begin{definition}[Latency utility]
\[
  U(t) =
  \begin{cases}
    1 & t \leq 300\,\text{ms} \\[4pt]
    0 & t \geq 10{,}000\,\text{ms} \\[4pt]
    1 - \dfrac{\ln(t/300)}{\ln(10000/300)} & \text{otherwise.}
  \end{cases}
\]
\end{definition}

Key landmark values:

\begin{table}[h!]
\centering
\caption{Utility at landmark TTFT values.}
\begin{tabular}{@{}rl@{}}
\toprule
TTFT $t$ & $U(t)$ \\
\midrule
$\leq 300\,\text{ms}$ & $1.000$ (perceptual immediacy) \\
$500\,\text{ms}$ & $0.854$ \\
$1{,}000\,\text{ms}$ & $0.657$ \\
$2{,}000\,\text{ms}$ & $0.459$ (attention degradation threshold) \\
$4{,}000\,\text{ms}$ & $0.261$ \\
$8{,}000\,\text{ms}$ & $0.064$ (abandonment risk threshold) \\
$10{,}000\,\text{ms}$ & $0.000$ (flow fully disrupted) \\
\bottomrule
\end{tabular}
\end{table}

\subsection{Session-level Metrics}

The model is applied \emph{per sample} $t_i$:
\[
  \bar{U} = \frac{1}{N}\sum_{i=1}^{N} U(t_i).
\]
Risk fractions:
\[
  f_{\text{deg}} = \frac{|\{i : t_i > 2000\}|}{N}, \qquad
  f_{\text{ab}} = \frac{|\{i : t_i > 8000\}|}{N}.
\]
The score returned to the composite pipeline is $\bar{U}$.

\subsection{Worked Example}
\label{sec:latency-utility-example}

We evaluate one realistic $N = 6$-turn session whose per-turn TTFT samples
are
\[
  (t_1, \ldots, t_6)
  = (250,\; 600,\; 1{,}500,\; 2{,}400,\; 3{,}000,\; 9{,}000)\ \text{ms}.
\]

\paragraph{Step 1 --- Normalising constant.}
The denominator of the log-linear branch depends only on the curve's two
endpoints, so we compute it once and reuse it for every sample.
\[
  \ln\!\left(\frac{10000}{300}\right) = \ln(33.333) = 3.5066.
\]

\paragraph{Step 2 --- Per-sample utilities.}
Each sample is pushed through $U(t)$ independently; we write out the second
sample in full and tabulate the remaining five identically.
\[
  U(t_2) = 1 - \frac{\ln(600/300)}{3.5066}
         = 1 - \frac{\ln 2}{3.5066}
         = 1 - \frac{0.6931}{3.5066}
         = 1 - 0.1977
         = 0.8023.
\]

\begin{table}[h!]
\centering
\caption{Per-sample utilities for the worked session.}
\label{tab:latency-per-sample}
\begin{tabular}{@{}crlrrr@{}}
\toprule
$i$ & $t_i$ & Branch of $U$ & $\ln(t_i/300)$ & $\ln(t_i/300)/3.5066$ & $U(t_i)$ \\
\midrule
1 & $250\,\text{ms}$   & $t \leq 300\,\text{ms}$ (clamp) & --- & --- & $1.0000$ \\
2 & $600\,\text{ms}$   & log-linear & $0.6931$ & $0.1977$ & $0.8023$ \\
3 & $1{,}500\,\text{ms}$ & log-linear & $1.6094$ & $0.4590$ & $0.5410$ \\
4 & $2{,}400\,\text{ms}$ & log-linear & $2.0794$ & $0.5930$ & $0.4070$ \\
5 & $3{,}000\,\text{ms}$ & log-linear & $2.3026$ & $0.6567$ & $0.3433$ \\
6 & $9{,}000\,\text{ms}$ & log-linear & $3.4012$ & $0.9700$ & $0.0300$ \\
\bottomrule
\end{tabular}
\end{table}

\paragraph{Step 3 --- Mean utility $\bar{U}$.}
The six per-sample utilities are averaged; this mean is the session-level
score, so every slow turn drags it down proportionally.
\[
  \sum_{i=1}^{6} U(t_i)
  = 1.0000 + 0.8023 + 0.5410 + 0.4070 + 0.3433 + 0.0300
  = 3.1236,
\]
\[
  \bar{U} = \frac{1}{N}\sum_{i=1}^{N} U(t_i)
          = \frac{3.1236}{6}
          = 0.5206.
\]

\paragraph{Step 4 --- Degradation fraction $f_{\text{deg}}$.}
We count the samples above the $2{,}000\,\text{ms}$ attention-degradation
threshold ($t_4$, $t_5$, $t_6$) and divide by $N$.
\[
  f_{\text{deg}} = \frac{|\{i : t_i > 2000\}|}{N}
                 = \frac{3}{6}
                 = 0.500 \quad (50.0\,\%).
\]

\paragraph{Step 5 --- Abandonment fraction $f_{\text{ab}}$.}
We count the samples above the $8{,}000\,\text{ms}$ abandonment-risk
threshold (only $t_6$) and divide by $N$.
\[
  f_{\text{ab}} = \frac{|\{i : t_i > 8000\}|}{N}
                = \frac{1}{6}
                = 0.167 \quad (16.7\,\%).
\]

The score returned to the composite pipeline is $\bar{U} = 0.5206$.  This
illustrates the motivation above: the session's median TTFT,
$(t_3 + t_4)/2 = 1{,}950\,\text{ms}$, sits below the degradation threshold
and would look acceptable as a point estimate, yet the per-sample mean
$\bar{U} = 0.52$ together with $f_{\text{ab}} = 1/6$ exposes that half the
turns degrade attention and one turn approaches full flow disruption.

\begin{table}[h!]
\centering
\caption{Summary of the worked latency-utility example.}
\label{tab:latency-worked}
\begin{tabular}{@{}llll@{}}
\toprule
Symbol & Meaning & Value (with unit) & What happens at this step \\
\midrule
$N$ & number of TTFT samples & 6 turns & raw input from telemetry \\
$t_i$ & per-turn TTFT samples & 250--9{,}000 ms & each pushed through $U$ \\
$\ln(10000/300)$ & normalising constant & 3.5066 & computed once; fixes slope \\
$U(t_i)$ & per-sample utilities & 1.0000--0.0300 & clamp at ends, log-linear between \\
$\bar{U}$ & mean utility & 0.5206 & $3.1236 / 6$; pipeline score \\
$f_{\text{deg}}$ & degradation fraction & 0.500 (50.0\,\%) & 3 of 6 samples $> 2{,}000$ ms \\
$f_{\text{ab}}$ & abandonment fraction & 0.167 (16.7\,\%) & 1 of 6 samples $> 8{,}000$ ms \\
\bottomrule
\end{tabular}
\end{table}

% ------------------------------------------------------------
\section{Token \& Cache Economics}
\label{sec:token-economics}
% ------------------------------------------------------------

\subsection{Pricing Model}

Prices are in USD per million tokens (MTok).  For each model $m$, let
$(p^m_{\text{in}},\, p^m_{\text{out}},\, p^m_{\text{cache}})$ be the
per-MTok costs for input, output, and cache-read tokens respectively.

Model resolution uses \emph{longest-prefix matching} so that
\texttt{gpt-4o-mini} is not consumed by the shorter key \texttt{gpt-4o}.

\subsection{Cost Formula}

For each model $m$ with token counts $(I_m, O_m, C_m)$:
\[
  \text{Cost}_m = \frac{I_m}{10^6}\,p^m_{\text{in}}
                + \frac{O_m}{10^6}\,p^m_{\text{out}}
                + \frac{C_m}{10^6}\,p^m_{\text{cache}},
  \qquad
  \text{Total} = \sum_m \text{Cost}_m.
\]

\subsection{Cache Savings}

Each cache-read token avoids paying the full input price:
\[
  \text{Savings}_m = \frac{C_m}{10^6}\,\max\!\left(0,\;p^m_{\text{in}} - p^m_{\text{cache}}\right).
\]

\subsection{Cache Hit Ratio (score)}

\begin{definition}[Cache hit ratio]
\[
  r_c = \frac{\sum_m C_m}{\sum_m (I_m + C_m)}.
\]
\end{definition}

A ratio $r_c \geq 0.5$ indicates that the prompt cache is doing heavy lifting;
$r_c < 0.1$ with large context ($> 50{,}000$ prompt tokens) flags fresh-context
waste.  The score fed to the efficiency component is $r_c$.

\subsection{Worked Example}
\label{sec:token-economics-example}

We cost one realistic session that used two models: a primary model
$m_1 = \texttt{claude-sonnet}$ for code generation and a cheaper
$m_2 = \texttt{gpt-4o-mini}$ for summarisation, with token counts
\[
  (I_{m_1}, O_{m_1}, C_{m_1}) = (120{,}000,\; 30{,}000,\; 480{,}000)\ \text{tokens},
\]
\[
  (I_{m_2}, O_{m_2}, C_{m_2}) = (40{,}000,\; 10{,}000,\; 60{,}000)\ \text{tokens}.
\]

\paragraph{Step 1 --- Price resolution.}
Each model name is resolved against the price sheet by longest-prefix
matching; note that \texttt{gpt-4o-mini} matches its own key rather than
being swallowed by the shorter prefix \texttt{gpt-4o}.
\[
  (p^{m_1}_{\text{in}},\, p^{m_1}_{\text{out}},\, p^{m_1}_{\text{cache}})
  = (3.00,\; 15.00,\; 0.30)\ \text{USD/MTok},
\]
\[
  (p^{m_2}_{\text{in}},\, p^{m_2}_{\text{out}},\, p^{m_2}_{\text{cache}})
  = (0.15,\; 0.60,\; 0.075)\ \text{USD/MTok}.
\]

\paragraph{Step 2 --- Cost of $m_1$.}
The three token counts are scaled to millions and multiplied by their
respective per-MTok prices, then summed.
\[
  \text{Cost}_{m_1}
  = \frac{120{,}000}{10^6} \times 3.00
  + \frac{30{,}000}{10^6} \times 15.00
  + \frac{480{,}000}{10^6} \times 0.30
\]
\[
  \phantom{\text{Cost}_{m_1}}
  = 0.12 \times 3.00 + 0.03 \times 15.00 + 0.48 \times 0.30
  = 0.3600 + 0.4500 + 0.1440
  = \$0.9540.
\]

\paragraph{Step 3 --- Cost of $m_2$.}
The same formula is applied to the summarisation model, whose prices are
an order of magnitude lower.
\[
  \text{Cost}_{m_2}
  = \frac{40{,}000}{10^6} \times 0.15
  + \frac{10{,}000}{10^6} \times 0.60
  + \frac{60{,}000}{10^6} \times 0.075
\]
\[
  \phantom{\text{Cost}_{m_2}}
  = 0.04 \times 0.15 + 0.01 \times 0.60 + 0.06 \times 0.075
  = 0.0060 + 0.0060 + 0.0045
  = \$0.0165.
\]

\paragraph{Step 4 --- Session total.}
The per-model costs are summed over the two models used.
\[
  \text{Total} = \sum_m \text{Cost}_m
               = 0.9540 + 0.0165
               = \$0.9705.
\]

\paragraph{Step 5 --- Cache savings.}
For each model, every cache-read token avoided paying the full input price,
so the saving is the cache volume times the price gap.
\[
  \text{Savings}_{m_1}
  = \frac{480{,}000}{10^6} \times \max(0,\; 3.00 - 0.30)
  = 0.48 \times 2.70
  = \$1.2960,
\]
\[
  \text{Savings}_{m_2}
  = \frac{60{,}000}{10^6} \times \max(0,\; 0.15 - 0.075)
  = 0.06 \times 0.075
  = \$0.0045,
\]
\[
  \text{Savings} = 1.2960 + 0.0045 = \$1.3005.
\]

\paragraph{Step 6 --- Cache hit ratio $r_c$.}
Finally we measure what fraction of all prompt tokens were served from
cache rather than paid at the full input price.
\[
  r_c = \frac{\sum_m C_m}{\sum_m (I_m + C_m)}
      = \frac{480{,}000 + 60{,}000}{(120{,}000 + 480{,}000) + (40{,}000 + 60{,}000)}
      = \frac{540{,}000}{700{,}000}
      = 0.771.
\]

Since $r_c = 0.771 \geq 0.5$, the prompt cache is doing heavy lifting.
Without caching the session would have cost
$0.9705 + 1.3005 = \$2.2710$, so cache reads cut spend by
$1.3005 / 2.2710 \approx 57\,\%$.  The score fed to the efficiency
component is $r_c = 0.771$.

\begin{table}[h!]
\centering
\caption{Summary of the worked token-economics example.}
\label{tab:token-economics-worked}
\begin{tabular}{@{}llll@{}}
\toprule
Symbol & Meaning & Value (with unit) & What happens at this step \\
\midrule
$(I_{m_1}, O_{m_1}, C_{m_1})$ & $m_1$ token counts & $(120, 30, 480)$ kTok & raw input from telemetry \\
$(I_{m_2}, O_{m_2}, C_{m_2})$ & $m_2$ token counts & $(40, 10, 60)$ kTok & raw input from telemetry \\
$p^{m_1}$ & $m_1$ prices (in/out/cache) & $(3.00, 15.00, 0.30)$ \$/MTok & prefix-matched price sheet \\
$p^{m_2}$ & $m_2$ prices (in/out/cache) & $(0.15, 0.60, 0.075)$ \$/MTok & longest prefix beats \texttt{gpt-4o} \\
$\text{Cost}_{m_1}$ & primary-model cost & \$0.9540 & three price terms summed \\
$\text{Cost}_{m_2}$ & summariser cost & \$0.0165 & same formula, cheaper sheet \\
Total & session cost & \$0.9705 & sum over both models \\
Savings & cache savings & \$1.3005 & cache volume $\times$ price gap \\
$r_c$ & cache hit ratio & 0.771 (77.1\,\%) & $540\text{k} / 700\text{k}$; efficiency score \\
\bottomrule
\end{tabular}
\end{table}

% ------------------------------------------------------------
\section{Lexical Frustration Classifier}
\label{sec:frustration}
% ------------------------------------------------------------

\subsection{Design Philosophy}

This is a \textbf{report-only} signal, deliberately excluded from the composite
quality score.  Lexicon heuristics produce false positives (``no worries'',
sarcasm, language bias) and false negatives (terse professional frustration).
Every flagged message carries its reasons so that a human can judge.  Inclusion
in the composite should follow validation against real sessions.

\subsection{Per-message Score}

Five signals are combined per user message $x_i$:

\begin{definition}[Message frustration score]
\[
  f_i = \min\!\left(1,\;
    0.45\,\min(h^s_i,2)
    + 0.15\,\min(h^m_i,3)
    + 0.30\,\mathbf{1}\!\left[J(x_i, x_{i-1}) \geq 0.6\right]
    + 0.15\,\mathbf{1}\!\left[\text{CAPS}(x_i)\right]
    + 0.10\,\mathbf{1}\!\left[\text{punct}(x_i)\right]
    + 0.20\,\mathbf{1}\!\left[\text{corr}(x_i)\right]
  \right)
\]
\end{definition}

where $h^s_i$ and $h^m_i$ are counts of strong and mild lexicon markers (bilingual
EN/PL), and the indicator predicates are:
\begin{itemize}
  \item $J(a,b) \geq 0.6$: Jaccard word-set similarity with previous message $\geq 0.6$
        (rephrasing / repair signal);
  \item $\text{CAPS}(x)$: $>60\%$ uppercase letters over $\geq 12$ letter characters;
  \item $\text{punct}(x)$: presence of \texttt{!!}, \texttt{??}, or \texttt{?!};
  \item $\text{corr}(x)$: $|x| < 20$ characters \emph{and} starts with
        ``no'', ``nie'', ``stop'', ``wrong'', or ``\'{z}le''.
\end{itemize}

\subsection{Jaccard Word-set Similarity}

Tokenize by splitting on whitespace and punctuation, keeping tokens of length $> 2$:

\begin{definition}[Jaccard similarity]
\[
  J(a, b) = \frac{|\calW_a \cap \calW_b|}{|\calW_a \cup \calW_b|}
           = \frac{|\calW_a \cap \calW_b|}{|\calW_a| + |\calW_b| - |\calW_a \cap \calW_b|}.
\]
\end{definition}

\subsection{Session Frustration Index}

The mean of per-message scores averages away isolated angry messages.
The index is pulled \emph{halfway toward the peak}:

\begin{definition}[Session frustration index]
\[
  \boxed{F = \bar{f} + 0.5\,\left(\max_i f_i - \bar{f}\right).}
\]
\end{definition}

\begin{table}[h!]
\centering
\caption{Frustration index interpretation.}
\begin{tabular}{@{}ll@{}}
\toprule
$F$ & Interpretation \\
\midrule
$\geq 0.5$ & Clear frustration signals --- review flagged messages \\
$\in [0.25, 0.5)$ & Mild friction in conversation tone \\
$< 0.25$ & No meaningful frustration signals \\
\bottomrule
\end{tabular}
\end{table}

\subsection{Worked Example}
\label{sec:frustration-example}

We classify one realistic 4-message exchange:
\begin{itemize}
  \item $x_1$: \textit{``Add pagination to the results table''}
  \item $x_2$: \textit{``add pagination to the results table component''}
  \item $x_3$: \textit{``this is STILL NOT WORKING??''}
  \item $x_4$: \textit{``no, revert that''}
\end{itemize}

\paragraph{Step 1 --- Baseline message $f_1$.}
The opening request carries no lexicon markers, has no previous message
to rephrase, and is typographically calm, so every term is zero.
\[
  f_1 = \min(1,\; 0) = 0.00.
\]

\paragraph{Step 2 --- Rephrasing signal $f_2$.}
Message $x_2$ nearly repeats $x_1$, so only the Jaccard indicator fires.
Tokenising with length $> 2$ gives
$\calW_{x_1} = \{\text{add}, \text{pagination}, \text{the},
\text{results}, \text{table}\}$ and
$\calW_{x_2} = \calW_{x_1} \cup \{\text{component}\}$:
\[
  J(x_2, x_1) = \frac{|\calW_{x_2} \cap \calW_{x_1}|}{|\calW_{x_2} \cup \calW_{x_1}|}
              = \frac{5}{6}
              = 0.833 \;\geq\; 0.6,
\]
\[
  f_2 = \min(1,\; 0.30 \times 1) = 0.30.
\]

\paragraph{Step 3 --- Peak message $f_3$.}
Message $x_3$ trips four signals: one strong marker
(\textit{``not working''}, $h^s_3 = 1$), one mild marker
(\textit{``still''}, $h^m_3 = 1$), sustained CAPS (15 of 21 letters
uppercase, $71\,\% > 60\,\%$ over $\geq 12$ letters), and a
\texttt{??}~punctuation burst.
\[
  f_3 = \min\!\left(1,\;
    0.45 \times \min(1,2)
    + 0.15 \times \min(1,3)
    + 0.15
    + 0.10
  \right)
\]
\[
  \phantom{f_3}
  = \min(1,\; 0.45 + 0.15 + 0.15 + 0.10)
  = \min(1,\; 0.85)
  = 0.85.
\]

\paragraph{Step 4 --- Short corrective reply $f_4$.}
Message $x_4$ contains two mild markers (\textit{``no,''} and
\textit{``revert''}, $h^m_4 = 2$) and is a short corrective: 15
characters $< 20$, starting with ``no''.
\[
  f_4 = \min\!\left(1,\; 0.15 \times \min(2,3) + 0.20\right)
      = \min(1,\; 0.30 + 0.20)
      = 0.50.
\]

\paragraph{Step 5 --- Session mean $\bar{f}$.}
The four per-message scores are averaged.
\[
  \bar{f} = \frac{0.00 + 0.30 + 0.85 + 0.50}{4}
          = \frac{1.65}{4}
          = 0.4125.
\]

\paragraph{Step 6 --- Session frustration index $F$.}
The mean alone would dilute the angry message, so the index is pulled
halfway toward the peak $\max_i f_i = 0.85$.
\[
  F = \bar{f} + 0.5\left(\max_i f_i - \bar{f}\right)
    = 0.4125 + 0.5 \times (0.85 - 0.4125)
    = 0.4125 + 0.2188
    = 0.631.
\]

Since $F = 0.631 \geq 0.5$, the session reports \emph{clear frustration
signals}; messages $x_3$ and $x_4$ are flagged with their reasons for
human review, per the report-only design philosophy.

\begin{table}[h!]
\centering
\caption{Summary of the worked frustration example.}
\label{tab:frustration-worked}
\begin{tabular}{@{}llll@{}}
\toprule
Symbol & Meaning & Value (with unit) & What happens at this step \\
\midrule
$f_1$ & baseline message score & 0.00 & no signals fire \\
$J(x_2, x_1)$ & rephrasing similarity & 0.833 & $5/6$ shared tokens; $\geq 0.6$ \\
$f_2$ & rephrased message score & 0.30 & Jaccard indicator only \\
$f_3$ & peak message score & 0.85 & strong + mild + CAPS + punct \\
$f_4$ & corrective reply score & 0.50 & two mild markers + short ``no'' \\
$\bar{f}$ & session mean & 0.4125 & $1.65 / 4$ \\
$\max_i f_i$ & session peak & 0.85 & message $x_3$ \\
$F$ & frustration index & 0.631 & mean pulled halfway to peak \\
\bottomrule
\end{tabular}
\end{table}

% ------------------------------------------------------------
\section{Order-Statistic Percentile}
\label{sec:percentile}
% ------------------------------------------------------------

All percentile computations in the pipeline use the \emph{nearest-rank}
(ceil) method.

\begin{definition}[Nearest-rank percentile]
Given a sorted list $v_{(1)} \leq \cdots \leq v_{(n)}$ and quantile
$p \in (0, 1]$:
\[
  \text{Percentile}(V, p) = v_{\!\left(\,\clamp{\lceil p\,n \rceil - 1}{0}{n-1}\right)}.
\]
\end{definition}

\begin{remark}
The nearest-rank method always returns an \emph{observed} value, not a
hypothetical interpolated one.  This is desirable for latency reporting
(a real TTFT measurement, not a weighted average of two measurements).
For small $n$ ($\leq 10$), divergence from linear interpolation can be
several hundred milliseconds at p95; the nearest-rank value is less
sensitive to outliers in that regime.
\end{remark}

\paragraph{Edge cases.}
Empty list returns~0.  A single-element list returns that element for all~$p$.

\subsection{Worked Example}
\label{sec:percentile-example}

We evaluate three quantiles over one session's $n = 8$ sorted TTFT
samples (the clamp maps the rank to a zero-based position in the sorted
list):
\[
  V = (420,\; 610,\; 780,\; 950,\; 1{,}200,\; 1{,}900,\; 2{,}600,\; 4{,}100)\ \text{ms}.
\]

\paragraph{Step 1 --- Median ($p = 0.50$).}
We scale the quantile by $n$, take the ceiling, shift to a zero-based
position, and clamp it into range.
\[
  \lceil 0.50 \times 8 \rceil - 1 = \lceil 4 \rceil - 1 = 3,
  \qquad
  \clamp{3}{0}{7} = 3,
\]
\[
  \text{Percentile}(V, 0.50) = v_{(3)} = 950\,\text{ms}.
\]

\paragraph{Step 2 --- Lower quartile ($p = 0.25$).}
The same three operations with $p = 0.25$.
\[
  \lceil 0.25 \times 8 \rceil - 1 = \lceil 2 \rceil - 1 = 1,
  \qquad
  \clamp{1}{0}{7} = 1,
\]
\[
  \text{Percentile}(V, 0.25) = v_{(1)} = 610\,\text{ms}.
\]

\paragraph{Step 3 --- Tail ($p = 0.95$).}
At the tail the ceiling rounds $7.6$ up to the last position, so the
method returns the largest observed sample.
\[
  \lceil 0.95 \times 8 \rceil - 1 = \lceil 7.6 \rceil - 1 = 8 - 1 = 7,
  \qquad
  \clamp{7}{0}{7} = 7,
\]
\[
  \text{Percentile}(V, 0.95) = v_{(7)} = 4{,}100\,\text{ms}.
\]

This illustrates the small-$n$ divergence noted above: linear
interpolation at $p = 0.95$ would report
$2{,}600 + 0.65 \times (4{,}100 - 2{,}600) = 3{,}575\,\text{ms}$, a
hypothetical value $525\,\text{ms}$ below the real measurement the
nearest-rank method returns.

\begin{table}[h!]
\centering
\caption{Summary of the worked percentile example.}
\label{tab:percentile-worked}
\begin{tabular}{@{}llll@{}}
\toprule
Symbol & Meaning & Value (with unit) & What happens at this step \\
\midrule
$n$ & sample count & 8 samples & sorted TTFT list \\
$p$ & quantiles evaluated & 0.25, 0.50, 0.95 & scaled by $n$, ceiled \\
$\text{Percentile}(V, 0.50)$ & median & 950 ms & position 3 after clamp \\
$\text{Percentile}(V, 0.25)$ & lower quartile & 610 ms & position 1 after clamp \\
$\text{Percentile}(V, 0.95)$ & tail latency & 4{,}100 ms & ceiling lands on last sample \\
--- & interpolated p95 & 3{,}575 ms & 525 ms below observed; not used \\
\bottomrule
\end{tabular}
\end{table}

% ------------------------------------------------------------
\section{Custom Protobuf Decoder}
\label{sec:proto}
% ------------------------------------------------------------

The collector ingests OTLP/proto payloads without the official
\texttt{Google.Protobuf} NuGet package, using a hand-written wire-format
parser.  This section documents the decoding rules for completeness.

\subsection{Wire Types}

The protobuf binary encoding maps each field to one of six wire types:

\begin{table}[h!]
\centering
\caption{Protobuf wire types used in OTLP payloads.}
\begin{tabular}{@{}rll@{}}
\toprule
Wire type & Name & Used for \\
\midrule
0 & Varint & field numbers, booleans, enum values \\
1 & 64-bit & \texttt{fixed64}, \texttt{sfixed64}, \texttt{double} \\
2 & Length-delimited & strings, bytes, embedded messages \\
5 & 32-bit & \texttt{fixed32}, \texttt{sfixed32}, \texttt{float} \\
\bottomrule
\end{tabular}
\end{table}

\subsection{Varint Decoding}

A varint encodes a 64-bit integer as a sequence of bytes; the MSB of
each byte signals continuation:
\[
  v = \sum_{k=0}^{K-1} (b_k \;\&\; 0x7F) \ll 7k,
  \quad \text{where } b_K \;\&\; 0x80 = 0.
\]

\subsection{Tag Decoding}

Each field tag encodes the field number and wire type:
\[
  \text{tag} = (\text{field\_number} \ll 3) \;|\; \text{wire\_type}.
\]

The \texttt{OtlpDecoder} traverses resource spans, scope spans, and
individual \texttt{Span} records, extracting attributes by name
(e.g.\ \texttt{gen\_ai.usage.input\_tokens}) and routing them to the
appropriate session aggregate.

\subsection{Worked Example}
\label{sec:proto-example}

We decode a realistic three-byte fragment carrying a token count of a
\texttt{gen\_ai.usage} attribute --- say field number~2 holding the value
150 tokens:
\[
  \texttt{0x10}\;\; \texttt{0x96}\;\; \texttt{0x01}.
\]

\paragraph{Step 1 --- Tag decoding.}
The first byte is the tag; splitting it per the tag formula yields the
field number (upper bits) and the wire type (lower three bits).
\[
  \text{tag} = \texttt{0x10} = 16,
\]
\[
  \text{field\_number} = 16 \gg 3 = 2,
  \qquad
  \text{wire\_type} = 16 \;\&\; 7 = 0 \quad (\text{varint}).
\]

\paragraph{Step 2 --- Continuation bits.}
Wire type 0 tells the decoder to read a varint next; each byte's MSB says
whether another byte follows.
\[
  b_0 = \texttt{0x96} = 1001\,0110_2:\quad
  b_0 \;\&\; \texttt{0x80} = \texttt{0x80} \neq 0
  \;\Rightarrow\; \text{continue},
\]
\[
  b_1 = \texttt{0x01} = 0000\,0001_2:\quad
  b_1 \;\&\; \texttt{0x80} = 0
  \;\Rightarrow\; \text{stop after } K = 2 \text{ bytes}.
\]

\paragraph{Step 3 --- Payload extraction.}
Masking each byte with \texttt{0x7F} keeps its 7 payload bits.
\[
  b_0 \;\&\; \texttt{0x7F} = \texttt{0x16} = 22,
  \qquad
  b_1 \;\&\; \texttt{0x7F} = \texttt{0x01} = 1.
\]

\paragraph{Step 4 --- Little-endian reassembly.}
The payloads are shifted into place, 7 bits per byte, least-significant
group first, and summed per the varint formula.
\[
  v = \sum_{k=0}^{1} (b_k \;\&\; \texttt{0x7F}) \ll 7k
    = 22 \ll 0 \;+\; 1 \ll 7
    = 22 + 128
    = 150.
\]

The decoder therefore records \texttt{input\_tokens} $= 150$ tokens ---
a two-byte varint, versus eight bytes for a fixed-width 64-bit integer.

\begin{table}[h!]
\centering
\caption{Summary of the worked protobuf-decoding example.}
\label{tab:proto-worked}
\begin{tabular}{@{}llll@{}}
\toprule
Symbol & Meaning & Value (with unit) & What happens at this step \\
\midrule
tag & first byte & \texttt{0x10} (16) & split into field and wire type \\
field\_number & which field & 2 & $16 \gg 3$ \\
wire\_type & how to read it & 0 (varint) & $16 \;\&\; 7$ \\
$b_0$ & first varint byte & \texttt{0x96} & MSB set: continue \\
$b_1$ & second varint byte & \texttt{0x01} & MSB clear: stop, $K = 2$ \\
$b_k \;\&\; \texttt{0x7F}$ & payload bits & 22, 1 & 7 bits kept per byte \\
$v$ & decoded value & 150 tokens & $22 + (1 \ll 7)$ \\
\bottomrule
\end{tabular}
\end{table}

% ------------------------------------------------------------
\section{Pipeline Integration}
\label{sec:pipeline}
% ------------------------------------------------------------

\subsection{InsightPipeline}

All analyzers implement \texttt{IInsightAnalyzer} and are composed in
\texttt{InsightPipeline}, which calls them sequentially and isolates
faults---a crash in one analyzer does not abort the others.  Results are
returned as \texttt{InsightReport} records:
\[
  \texttt{InsightReport} = (\text{title},\, \text{algorithm},\, \text{status},\,
    \text{score} \in [0,1] \cup \{\text{null}\},\, \text{metrics},\, \text{findings}).
\]

\subsection{Session Kind Classification}

Before quality scoring, sessions are classified by scanning prompt prefixes:

\begin{table}[h!]
\centering
\caption{Session kind classification rules.}
\begin{tabular}{@{}ll@{}}
\toprule
Kind & Detection rule \\
\midrule
\texttt{UserChat} & No internal prefix detected \\
\texttt{InternalTitleGeneration} & Prompt starts with title-generation prefix \\
\texttt{InternalSummary} & Prompt starts with summary prefix \\
\texttt{InternalHelper} & Prompt starts with helper prefix \\
\texttt{Unattributed} & No transcript available \\
\bottomrule
\end{tabular}
\end{table}

Internal sessions (title generation, summarisation) are excluded from
quality aggregation to avoid inflating scores with trivially perfect
short sessions.

% ------------------------------------------------------------
\section{Discussion}
\label{sec:discussion}
% ------------------------------------------------------------

\subsection{Design Choices and Trade-offs}

\paragraph{Weight renormalisation (v2).}
The v1 composite engine contributed a fixed neutral prior~0.7 at full weight
for absent components, compressing all ordinary sessions toward 79--80 and
destroying discrimination.  The v2 renormalisation scheme allows the score to
reflect only what is actually observed.  The cost is that a session with only
reliability data has a score of $v_R \times 100$, which can be misleadingly
high or low---hence the confidence signal is essential.

\paragraph{Self-referential latency in TFRA.}
Using session-median TTFT as the reference for latency penalties means that a
slow model is not punished for being slow---only for \emph{inconsistency}
within its own distribution.  This is intentional: a developer who chooses
Claude Opus accepts higher latency; the friction signal should flag unexpected
spikes, not the baseline.

\paragraph{Frustration as report-only.}
Lexicon heuristics are noisy.  Promoting the frustration index into the
composite score before validating it against real labelled sessions would
introduce a signal with unknown false-positive rate into a number that drives
dashboard decisions.  The report-only design follows the principle of doing no
harm until a signal is validated.

\subsection{Upgrade Paths}

\begin{itemize}
  \item \textbf{Frustration}: replace the lexicon heuristic with a SPUR-style
        learned rubric or a lightweight local model trained on session transcripts.
  \item \textbf{Throughput}: add a semantic similarity component using
        diff-level embeddings to detect intent-preserving refactors.
  \item \textbf{Latency}: incorporate model-specific SLA targets so that
        latency utility curves are calibrated per model rather than universal.
  \item \textbf{Composite}: learn component weights from labelled ``good session''
        datasets rather than using hand-tuned values.
\end{itemize}

% ------------------------------------------------------------
\section{Conclusion}
\label{sec:conclusion}
% ------------------------------------------------------------

This paper has formally defined the eight algorithms that constitute the
CopilotScope quality measurement pipeline.  Together they provide a
\emph{multi-dimensional, telemetry-only, confidence-aware} view of AI
coding-assistant session quality---without requiring prompt inspection,
LLM judges, or external API calls at scoring time.

The framework is implemented in C\# (.NET~8) and is open for extension:
new analyzers implement \texttt{IInsightAnalyzer} and register with the
pipeline in a single line of code.  The Python research notebooks in the
companion \texttt{research/} directory reproduce every formula and numerical
example in this paper and provide visualisations suitable for further
investigation.

% ------------------------------------------------------------
\appendix
% ------------------------------------------------------------

\section{Symbol Table}
\label{app:symbols}

\begin{table}[h!]
\centering
\begin{tabular}{@{}ll@{}}
\toprule
Symbol & Meaning \\
\midrule
$Q$ & Composite quality score $\in [0, 100]$ \\
$C$ & Weight coverage of informative components \\
$\calI$ & Set of informative component indices \\
$w_i$ & Base weight of component $i$ \\
$v_i$ & Value of component $i \in [0,1]$ \\
$n_i$ & Sample count of component $i$ \\
$\rho_i$ & Sample ramp of component $i = \min(1, n_i/5)$ \\
$\hat{e}$ & Weighted error rate (reliability) \\
$\alpha$ & Edit acceptance ratio \\
$\bar{s}$ & Mean survival score \\
$\bar{\tau}$ & Mean TFRA turn score (friction) \\
$\tau_i$ & TFRA score for turn $i$ \\
$m$ & Session median TTFT \\
$U(t)$ & Latency utility at TTFT $t$ \\
$\bar{U}$ & Mean per-sample latency utility \\
$f_{\text{deg}}$ & Fraction of responses $> 2\,\text{s}$ \\
$f_{\text{ab}}$ & Fraction of responses $> 8\,\text{s}$ \\
$r_c$ & Prompt-cache hit ratio \\
$f_i$ & Per-message frustration score \\
$F$ & Session frustration index \\
$J(a,b)$ & Jaccard word-set similarity \\
\bottomrule
\end{tabular}
\end{table}

\section{Base Weights Sensitivity}
\label{app:sensitivity}

The base weights $(0.25, 0.20, 0.20, 0.15, 0.10, 0.10)$ sum to $1.0$.
The table below shows how the composite score changes under uniform
$\pm 0.05$ perturbation to the reliability weight, holding all others fixed
and renormalising, for the ``perfect session'' scenario ($v_i = 1$ for all $i$).
Since all values are~1, the score is always~100 regardless of weight choice
(the formula reduces to $100 \cdot C / C = 100$).  The sensitivity only
manifests when components have differing values---which is the practical case.

\end{document}
